\textbf{\textit{Resumen}} - El surgimiento de ecosistemas de gobernanza descentralizada y registros académicos inmutables ha posicionado a la tecnología Blockchain como un pilar para la integridad de datos. Sin embargo, la complejidad técnica de arquitecturas como Polkadot, basadas en el framework Substrate, introduce desafíos significativos en el aseguramiento de la calidad (QA) y la detección de vulnerabilidades en Smart Contracts. Este trabajo de investigación analiza el estado actual del Testing en entornos Blockchain, enfocándose específicamente en la verificación y validación de sistemas de gestión de notas académicas sobre la red Polkadot. Se exploran técnicas modernas de Continuous Testing, pruebas de mutación y análisis formal aplicados a entornos descentralizados. El estudio concluye que, si bien las herramientas de automatización actuales están evolucionando, la naturaleza inmutable de la cadena de bloques exige paradigmas de testing más rigurosos antes de la implantación en entornos de producción.
\vspace{0.5em}

\textbf{\textit{Index Terms}} - Blockchain, Polkadot, Sistema de notas, Seguridad, DApp, Descentralizado.